% --- Archon physics formalization source begin ---
% archon:physics
% archon:covers IPhO2026Problems/problem_IPhO_2026_1_C_1.lean
% archon:source-report reports/ipho_2026/problem_IPhO_2026_1_C_1.source.json
% archon:problem-id IPhO_2026_1
% archon:part-id C.1

\chapter{Physics problem IPhO\_2026\_1\_C\_1}
\label{ch:IPhO2026Problems_problem_IPhO_2026_1_C_1}

\paragraph{Problem source.}
A photon of angular frequency omega is absorbed by an ozone molecule O3 at rest,
dissociating it into O2 and O.  Let U\_i and U\_f be the ground-state energies of
O3 and O2 and define Delta U = U\_f - U\_i.  The outgoing O2 momentum makes angle
theta with the incident photon.  Treat the oxygen fragments classically and
non-relativistically, take the mass of an oxygen atom to be m, and use photon
momentum p\_gamma = E\_gamma/c = hbar*omega/c.

Current subquestion:
Determine the minimum angular frequency omega\_min required for dissociation at outgoing O2 angle theta, in terms of hbar, c, theta, Delta U, and m.

\paragraph{Current subquestion.}
Determine the minimum angular frequency omega\_min required for dissociation at outgoing O2 angle theta, in terms of hbar, c, theta, Delta U, and m.

\paragraph{Recorded answer/context.}
For theta <= pi/2, omega\_min = 3*m*c\textasciicircum{}2*[1 - sqrt(1 - (2*Delta U/(3*m*c\textasciicircum{}2))*(1 + 2*sin(theta)\textasciicircum{}2))]/[hbar*(1 + 2*sin(theta)\textasciicircum{}2)]. For theta >= pi/2 use the same threshold evaluated at theta = pi/2.

\paragraph{Figure/image path.}
/root/proposal\_for\_physic/science-mango/ipho\_2026\_source/image/T1\_page-3.png

\paragraph{Formalization target.}
create a compiling Lean file with sorry bodies at `IPhO2026Problems/problem\_IPhO\_2026\_1\_C\_1.lean`. The Lean declarations must preserve the physical quantities, dimensions or dimensional roles, figure labels, governing-law hypotheses, and final relation expressed by this problem.
Use Mathlib/Physlib names found through LeanExplore where available. If a physics API is missing, introduce faithful local abstractions rather than scalar placeholder aliases.

\paragraph{Iteration 004 proof-review redraft contract.}
The scalar event-balance lemma must assume the applicability conditions on the
physical parameters.  In particular, positivity of \(\hbar\), the photon
frequency, and \(c\) makes the quantity called the photon-momentum magnitude
nonnegative.  Do not assert the radial formula for arbitrary signed parameter
readouts.  Preserve the already established threshold expression and the
downstream theorem statements.

\begin{theorem}[Physics formalization target]
  \label{thm:physics:IPhO_2026_1_C_1:target}
  \lean{IPhO2026Problem1C1.minimumAngularFrequency_isDissociationThreshold}
  \uses{def:physics:IPhO_2026_1_C_1:aux001, def:physics:IPhO_2026_1_C_1:aux002, def:physics:IPhO_2026_1_C_1:aux003, def:physics:IPhO_2026_1_C_1:aux004, def:physics:IPhO_2026_1_C_1:aux005, def:physics:IPhO_2026_1_C_1:aux006, def:physics:IPhO_2026_1_C_1:aux007, def:physics:IPhO_2026_1_C_1:aux008, def:physics:IPhO_2026_1_C_1:aux009, def:physics:IPhO_2026_1_C_1:aux010, def:physics:IPhO_2026_1_C_1:aux011, def:physics:IPhO_2026_1_C_1:aux012, def:physics:IPhO_2026_1_C_1:aux013, def:physics:IPhO_2026_1_C_1:aux014, def:physics:IPhO_2026_1_C_1:aux015, lem:physics:IPhO_2026_1_C_1:aux016, lem:physics:IPhO_2026_1_C_1:aux017, lem:physics:IPhO_2026_1_C_1:aux018, def:physics:IPhO_2026_1_C_1:aux019, def:physics:IPhO_2026_1_C_1:aux020, lem:physics:IPhO_2026_1_C_1:aux021, lem:physics:IPhO_2026_1_C_1:aux022}
  The explicit formula itself is proved to be the minimum required angular
  frequency; no threshold formula is assumed as a governing law.
\end{theorem}
\begin{proof}
  Use the typed geometry, governing-law, branch, and measurement interfaces listed in the declaration index.  Eliminate the intermediate readouts in their physical order and then evaluate the requested exact or uncertainty relation.
\end{proof}
\section*{Formal declaration index}
The following blocks record the typed carriers and bridge declarations used by the target.  Their dependency lists follow the declaration-level model in the covered Lean file.

\begin{definition}[Declaration MomentumPlane]
  \label{def:physics:IPhO_2026_1_C_1:aux001}
  \lean{IPhO2026Problem1C1.MomentumPlane}
  The two-dimensional momentum-coordinate space used for Figure 1c.
\end{definition}
\begin{proof}
  This is the typed definition of the stated physical carrier; its fields or defining equation provide the claimed data.
\end{proof}

\begin{definition}[Declaration figure1cIncidentDirection]
  \label{def:physics:IPhO_2026_1_C_1:aux002}
  \lean{IPhO2026Problem1C1.figure1cIncidentDirection}
  \uses{def:physics:IPhO_2026_1_C_1:aux001}
  The incident-photon direction (the horizontal dashed ray in Figure 1c).
\end{definition}
\begin{proof}
  This is the typed definition of the stated physical carrier; its fields or defining equation provide the claimed data.
\end{proof}

\begin{definition}[Declaration Parameters]
  \label{def:physics:IPhO_2026_1_C_1:aux003}
  \lean{IPhO2026Problem1C1.Parameters}
  Scalar readouts for the photodissociation experiment in one coherent unit system. The fields have dimensions energy · time, mass, speed, and energy, respectively. The oxygen atom's ground-state energy is the chosen zero of energy, as in the source energy balance.
\end{definition}
\begin{proof}
  This is the typed definition of the stated physical carrier; its fields or defining equation provide the claimed data.
\end{proof}

\begin{definition}[Declaration Parameters.energyGap]
  \label{def:physics:IPhO_2026_1_C_1:aux004}
  \lean{IPhO2026Problem1C1.Parameters.energyGap}
  \uses{def:physics:IPhO_2026_1_C_1:aux003}
  The dissociation energy gap ΔU = U\_f - U\_i.
\end{definition}
\begin{proof}
  This is the typed definition of the stated physical carrier; its fields or defining equation provide the claimed data.
\end{proof}

\begin{definition}[Declaration Parameters.Valid]
  \label{def:physics:IPhO_2026_1_C_1:aux005}
  \lean{IPhO2026Problem1C1.Parameters.Valid}
  \uses{def:physics:IPhO_2026_1_C_1:aux003, def:physics:IPhO_2026_1_C_1:aux004}
  Applicability conditions for the classical threshold formula. The last inequality keeps the relevant discriminants strictly positive and is the quantitative small-energy condition appropriate to the nonrelativistic model.
\end{definition}
\begin{proof}
  This is the typed definition of the stated physical carrier; its fields or defining equation provide the claimed data.
\end{proof}

\begin{definition}[Declaration photonEnergy]
  \label{def:physics:IPhO_2026_1_C_1:aux006}
  \lean{IPhO2026Problem1C1.photonEnergy}
  \uses{def:physics:IPhO_2026_1_C_1:aux003}
  Photon energy E\_γ = ℏω in the chosen coherent units.
\end{definition}
\begin{proof}
  This is the typed definition of the stated physical carrier; its fields or defining equation provide the claimed data.
\end{proof}

\begin{definition}[Declaration photonMomentumMagnitude]
  \label{def:physics:IPhO_2026_1_C_1:aux007}
  \lean{IPhO2026Problem1C1.photonMomentumMagnitude}
  \uses{def:physics:IPhO_2026_1_C_1:aux003, def:physics:IPhO_2026_1_C_1:aux006}
  Photon momentum magnitude p\_γ = E\_γ / c = ℏω/c.
\end{definition}
\begin{proof}
  This is the typed definition of the stated physical carrier; its fields or defining equation provide the claimed data.
\end{proof}

\begin{definition}[Declaration fragmentKineticEnergy]
  \label{def:physics:IPhO_2026_1_C_1:aux008}
  \lean{IPhO2026Problem1C1.fragmentKineticEnergy}
  \uses{def:physics:IPhO_2026_1_C_1:aux001, def:physics:IPhO_2026_1_C_1:aux003}
  Classical, nonrelativistic kinetic energy of the outgoing fragments. The oxygen molecule has mass 2m, hence its term is |p|²/(4m); the oxygen atom has mass m, hence its term is |q|²/(2m).
\end{definition}
\begin{proof}
  This is the typed definition of the stated physical carrier; its fields or defining equation provide the claimed data.
\end{proof}

\begin{definition}[Declaration DissociationEvent]
  \label{def:physics:IPhO_2026_1_C_1:aux009}
  \lean{IPhO2026Problem1C1.DissociationEvent}
  \uses{def:physics:IPhO_2026_1_C_1:aux001, def:physics:IPhO_2026_1_C_1:aux002, def:physics:IPhO_2026_1_C_1:aux003, def:physics:IPhO_2026_1_C_1:aux006, def:physics:IPhO_2026_1_C_1:aux007, def:physics:IPhO_2026_1_C_1:aux008}
  A single event satisfying the governing laws of the isolated photodissociation process. The named momentum fields distinguish the incoming photon from both outgoing fragments. The angle is Mathlib's unoriented angle in [0, π], matching the angle label θ in Figure 1c.
\end{definition}
\begin{proof}
  This is the typed definition of the stated physical carrier; its fields or defining equation provide the claimed data.
\end{proof}

\begin{definition}[Declaration KinematicallyAllowed]
  \label{def:physics:IPhO_2026_1_C_1:aux010}
  \lean{IPhO2026Problem1C1.KinematicallyAllowed}
  \uses{def:physics:IPhO_2026_1_C_1:aux003, def:physics:IPhO_2026_1_C_1:aux009}
  Dissociation at the stated angle and frequency is kinematically possible.
\end{definition}
\begin{proof}
  This is the typed definition of the stated physical carrier; its fields or defining equation provide the claimed data.
\end{proof}

\begin{definition}[Declaration radialFragmentKineticEnergy]
  \label{def:physics:IPhO_2026_1_C_1:aux011}
  \lean{IPhO2026Problem1C1.radialFragmentKineticEnergy}
  \uses{def:physics:IPhO_2026_1_C_1:aux003, def:physics:IPhO_2026_1_C_1:aux007}
  The radial form of the total fragment kinetic energy after eliminating the oxygen-atom momentum with momentum conservation. Here r = |p\_\{O₂\}|.
\end{definition}
\begin{proof}
  This is the typed definition of the stated physical carrier; its fields or defining equation provide the claimed data.
\end{proof}

\begin{definition}[Declaration angularFactor]
  \label{def:physics:IPhO_2026_1_C_1:aux012}
  \lean{IPhO2026Problem1C1.angularFactor}
  The angular factor 1 + 2 sin² θ appearing in the threshold.
\end{definition}
\begin{proof}
  This is the typed definition of the stated physical carrier; its fields or defining equation provide the claimed data.
\end{proof}

\begin{definition}[Declaration minimumFragmentKineticEnergy]
  \label{def:physics:IPhO_2026_1_C_1:aux013}
  \lean{IPhO2026Problem1C1.minimumFragmentKineticEnergy}
  \uses{def:physics:IPhO_2026_1_C_1:aux003, def:physics:IPhO_2026_1_C_1:aux007, def:physics:IPhO_2026_1_C_1:aux012}
  Infimum of the radial fragment kinetic energy at fixed photon momentum. For forward angles the minimizing oxygen-molecule momentum is nonnegative. For angles at or beyond π/2, the constrained infimum occurs as that momentum tends to zero.
\end{definition}
\begin{proof}
  This is the typed definition of the stated physical carrier; its fields or defining equation provide the claimed data.
\end{proof}

\begin{definition}[Declaration HasEnoughPhotonEnergy]
  \label{def:physics:IPhO_2026_1_C_1:aux014}
  \lean{IPhO2026Problem1C1.HasEnoughPhotonEnergy}
  \uses{def:physics:IPhO_2026_1_C_1:aux003, def:physics:IPhO_2026_1_C_1:aux004, def:physics:IPhO_2026_1_C_1:aux006, def:physics:IPhO_2026_1_C_1:aux013}
  The exact scalar feasibility condition obtained by minimizing over the outgoing oxygen-molecule momentum magnitude. At and beyond π/2 the infimum requires zero O₂ momentum, for which the angle would not be physically defined; consequently actual events satisfy a strict inequality there.
\end{definition}
\begin{proof}
  This is the typed definition of the stated physical carrier; its fields or defining equation provide the claimed data.
\end{proof}

\begin{definition}[Declaration IsDissociationThreshold]
  \label{def:physics:IPhO_2026_1_C_1:aux015}
  \lean{IPhO2026Problem1C1.IsDissociationThreshold}
  \uses{def:physics:IPhO_2026_1_C_1:aux003, def:physics:IPhO_2026_1_C_1:aux010}
  An infimum characterization of the minimum required angular frequency. This permits the backscattering threshold to be a limiting value even when no event with nonzero outgoing O₂ momentum occurs exactly at that value.
\end{definition}
\begin{proof}
  This is the typed definition of the stated physical carrier; its fields or defining equation provide the claimed data.
\end{proof}

\begin{lemma}[Declaration event\_scalar\_energy\_balance]
  \label{lem:physics:IPhO_2026_1_C_1:aux016}
  \lean{IPhO2026Problem1C1.event_scalar_energy_balance}
  \uses{def:physics:IPhO_2026_1_C_1:aux003, def:physics:IPhO_2026_1_C_1:aux004, def:physics:IPhO_2026_1_C_1:aux005, def:physics:IPhO_2026_1_C_1:aux006, def:physics:IPhO_2026_1_C_1:aux007, def:physics:IPhO_2026_1_C_1:aux008, def:physics:IPhO_2026_1_C_1:aux009, def:physics:IPhO_2026_1_C_1:aux011}
  Under the parameter-validity conditions, vector conservation and the
  Figure 1c angle imply the scalar energy law.
\end{lemma}
\begin{proof}
  Energy conservation first gives photon energy as the dissociation gap plus
  the two fragment kinetic energies.  Momentum conservation replaces the
  atomic momentum by photon momentum minus molecular momentum.  Validity and
  the event's positive frequency make \(\hbar\omega/c\) positive, so the norm
  of the photon vector is exactly this magnitude rather than its absolute
  value.  The Figure 1c angle identity then evaluates the cross term as
  \(2p_\gamma\lVert p_{O_2}\rVert\cos\theta\), yielding the stated radial
  expression.
\end{proof}

\begin{lemma}[Declaration radialFragmentKineticEnergy\_lower\_bound]
  \label{lem:physics:IPhO_2026_1_C_1:aux017}
  \lean{IPhO2026Problem1C1.radialFragmentKineticEnergy_lower_bound}
  \uses{def:physics:IPhO_2026_1_C_1:aux003, def:physics:IPhO_2026_1_C_1:aux005, def:physics:IPhO_2026_1_C_1:aux011, def:physics:IPhO_2026_1_C_1:aux013, lem:physics:IPhO_2026_1_C_1:aux016}
  The constrained radial kinetic energy is bounded below by the appropriate forward/backscattering infimum.
\end{lemma}
\begin{proof}
  Substitute the cited governing relations in order and use the stated positivity, orientation, and branch conditions to obtain the conclusion.
\end{proof}

\begin{lemma}[Declaration kinematicallyAllowed\_iff\_hasEnoughPhotonEnergy]
  \label{lem:physics:IPhO_2026_1_C_1:aux018}
  \lean{IPhO2026Problem1C1.kinematicallyAllowed_iff_hasEnoughPhotonEnergy}
  \uses{def:physics:IPhO_2026_1_C_1:aux003, def:physics:IPhO_2026_1_C_1:aux005, def:physics:IPhO_2026_1_C_1:aux010, def:physics:IPhO_2026_1_C_1:aux014, lem:physics:IPhO_2026_1_C_1:aux016, lem:physics:IPhO_2026_1_C_1:aux017}
  The vector event interface is neither opaque nor underdetermined: eliminating the two momentum vectors gives exactly the scalar minimized-energy condition.
\end{lemma}
\begin{proof}
  Substitute the cited governing relations in order and use the stated positivity, orientation, and branch conditions to obtain the conclusion.
\end{proof}

\begin{definition}[Declaration effectiveThresholdAngle]
  \label{def:physics:IPhO_2026_1_C_1:aux019}
  \lean{IPhO2026Problem1C1.effectiveThresholdAngle}
  Clamp the Figure 1c angle to the forward limiting value π/2.
\end{definition}
\begin{proof}
  This is the typed definition of the stated physical carrier; its fields or defining equation provide the claimed data.
\end{proof}

\begin{definition}[Declaration minimumAngularFrequencyExpression]
  \label{def:physics:IPhO_2026_1_C_1:aux020}
  \lean{IPhO2026Problem1C1.minimumAngularFrequencyExpression}
  \uses{def:physics:IPhO_2026_1_C_1:aux003, def:physics:IPhO_2026_1_C_1:aux004, def:physics:IPhO_2026_1_C_1:aux012, def:physics:IPhO_2026_1_C_1:aux019}
  The threshold expression obtained from the lower root of the energy quadratic. The factor 2 multiplying ΔU under the square root is present in the official solution and is also required by the numerical result in part C.2. The generated blueprint's recorded-answer line omitted this factor.
\end{definition}
\begin{proof}
  This is the typed definition of the stated physical carrier; its fields or defining equation provide the claimed data.
\end{proof}

\begin{lemma}[Declaration minimumAngularFrequencyExpression\_energy\_boundary]
  \label{lem:physics:IPhO_2026_1_C_1:aux021}
  \lean{IPhO2026Problem1C1.minimumAngularFrequencyExpression_energy_boundary}
  \uses{def:physics:IPhO_2026_1_C_1:aux003, def:physics:IPhO_2026_1_C_1:aux004, def:physics:IPhO_2026_1_C_1:aux005, def:physics:IPhO_2026_1_C_1:aux006, def:physics:IPhO_2026_1_C_1:aux013, def:physics:IPhO_2026_1_C_1:aux020, lem:physics:IPhO_2026_1_C_1:aux016, lem:physics:IPhO_2026_1_C_1:aux017, lem:physics:IPhO_2026_1_C_1:aux018}
  The displayed expression is the lower root of the minimized energy equation.
\end{lemma}
\begin{proof}
  Substitute the cited governing relations in order and use the stated positivity, orientation, and branch conditions to obtain the conclusion.
\end{proof}

\begin{lemma}[Declaration isDissociationThreshold\_unique]
  \label{lem:physics:IPhO_2026_1_C_1:aux022}
  \lean{IPhO2026Problem1C1.isDissociationThreshold_unique}
  \uses{def:physics:IPhO_2026_1_C_1:aux003, def:physics:IPhO_2026_1_C_1:aux015, lem:physics:IPhO_2026_1_C_1:aux016, lem:physics:IPhO_2026_1_C_1:aux017, lem:physics:IPhO_2026_1_C_1:aux018, lem:physics:IPhO_2026_1_C_1:aux021}
  An infimum threshold, when it exists, is unique.
\end{lemma}
\begin{proof}
  Substitute the cited governing relations in order and use the stated positivity, orientation, and branch conditions to obtain the conclusion.
\end{proof}

\begin{lemma}[Declaration minimumAngularFrequency\_eq]
  \label{lem:physics:IPhO_2026_1_C_1:aux024}
  \lean{IPhO2026Problem1C1.minimumAngularFrequency_eq}
  \uses{def:physics:IPhO_2026_1_C_1:aux003, def:physics:IPhO_2026_1_C_1:aux005, def:physics:IPhO_2026_1_C_1:aux015, def:physics:IPhO_2026_1_C_1:aux020, lem:physics:IPhO_2026_1_C_1:aux016, lem:physics:IPhO_2026_1_C_1:aux017, lem:physics:IPhO_2026_1_C_1:aux018, lem:physics:IPhO_2026_1_C_1:aux021, lem:physics:IPhO_2026_1_C_1:aux022, thm:physics:IPhO_2026_1_C_1:target}
  Equivalent value form: any scalar already identified semantically as the minimum dissociation frequency equals the explicit expression.
\end{lemma}
\begin{proof}
  Substitute the cited governing relations in order and use the stated positivity, orientation, and branch conditions to obtain the conclusion.
\end{proof}
% --- Archon physics formalization source end ---
